\documentclass[
  spinewidth=39mm,
  coverheight=190mm,
  coverwidth=125mm,
  bleedwidth=3mm,
  marklength=0mm
]{bookcover}

\usepackage[T1]{fontenc}
\usepackage[sfdefault]{AlegreyaSans}
\usepackage{xcolor}

\definecolor{stars}{HTML}{02174c}
\newcommand{\booktitle}{Die Weisheit des Brahmanen}
\newcommand{\bookauthor}{Friedrich Rückert}
\newcommand{\booksubtitle}{Ein Lehrgedicht in Bruchstücken}

\newbookcovercomponenttype{center rotate}{%
  \parbox[t][\partheight][c]{\partwidth}{%
    \begin{center}
       \rotatebox[origin=c]{90}{#1}
    \end{center}}}

\begin{document}
  \begin{bookcover}
    \bookcovercomponent{color}{bg whole}{black!90!white}
    \bookcovercomponent{color}{bg spine}{black!90!white}
    \bookcovercomponent{color}{bg front}{stars!70!black}
    \bookcovercomponent{normal}{front}[0.06\textwidth,0.06\textwidth,0.06\textwidth,0.06\textwidth]{
      \Huge\color{white}
      {\Large\bookauthor{}\\[0.5\baselineskip]}
      {\Huge\bfseries Die Weisheit \\ des Brahmanen\\[0.5\baselineskip]}
      {\huge\booksubtitle{}}\vfill
      \par
      PTOROS}


  \bookcovercomponent{normal}{back}{%
    \centering\large\color{white}\vspace{2em}
    \parbox{0.38\textwidth}{%
      \textbf{Friedrich Rückert}s imposantes Lehrgedicht \textit{Die Weisheit des Brahmanen}, zwischen 1836 und 1839 in sechs Bänden erschienen, bündelt in einzigartiger Weise Anregungen aus Christentum, Islam, Hinduismus und abendländischer Philosophie.

      \medskip

      \begin{quote}
        \emph{%
      Ein indischer Brahman, geboren auf der Flur,\\
      Der nichts gelesen als den Weda der Natur;\\
      Hat viel gesehn, gedacht, noch mehr geahnt, gefühlt,\\
      Und mit Betrachtungen die Leidenschaft gekühlt;\\
      Spricht bald was klar ihm ward, bald um sich’s klar zu machen,\\
      Von ihn angeh’nden halb, halb nicht angeh’nden Sachen.\\
      Er hat die Eigenheit, nur Einzelnes zu sehn,\\
      Doch alles Einzelne als Ganzes zu verstehn.\\
      Woran er immer nur sieht schimmern einen Glanz,\\
      Wird ein Betkügelchen an seinem Rosenkranz.
        }\\
        \mbox{}\hfill(I. 1.)
      \end{quote}
      \medskip
      \begin{quote}
        \emph{%
          Gar manches sagt nicht rein brahman’sches der Brahman;\\
          Sei es rein menschlich nur, so nehmen wir es an.\\
          Doch dieses, was aus gar zu fremden Augen schaut,\\
          Hat ein europischer Bekannter ihm vertraut.
        }\\
        \mbox{}\hfill(XI. 1.)
      \end{quote}
      \medskip
      \begin{quote}
        \emph{%
          Ich glaube nicht, daß ich viel eignes neues lehre,\\
          Noch durch mein Scherflein Witz den Schatz der Weisheit mehre.\\
          Doch denk’ ich von der Müh mir zweierlei Gewinn;\\
          Einmal, daß ich nun selbst an Einsicht weiter bin;\\
          Sodann, daß doch dadurch an manchen Mann wird kommen\\
          Manches, wovon er sonst gar hätte nichts vernommen.\\
          Und auch der dritte Grund scheint werth nicht des Gelächters:\\
          Daß, wer dis Büchlein list, derweil doch list kein schlechters.
        }\\
        \mbox{}\hfill(VIII. 1.)
      \end{quote}
      \bigskip
    }}

    \bookcovercomponent{center rotate}{spine}{
      \color{white}\Large
      \bookauthor{} \hspace{8.4cm} \textbf{\booktitle{}}
    }
  \end{bookcover}
\end{document}
